% UNREGISTERED D.57 candidate. Existing main sections are frozen.
\section*{Draft D.57: every nonzero D.20 class fails at a finite derived stage}

\noindent\textbf{Candidate statement.}
Let $\ell\in\{3,5\}$, $n=\ell^r$, and let
$s_a$ be the actual full metabelian Fermat family of D.20,
parametrized by $a\in\mathbb Z_\ell^{d_n}$ with
$d_n=(n^2-1)(n^2+3)/4$.  Then:
\begin{enumerate}
\item The restrictions of the $s_a$ to $G_{\mathbb Q_\ell}$
are pairwise nonconjugate in the full metabelian group.  For
$a\ne0$, the local section is neither point nor cusp induced.
\item No $a\ne0$ has a lift to the plain geometrically
prosolvable quotient, even locally at $\ell$.  This conclusion
holds for both the three-power and five-power towers.
\item For each $a\ne0$ there is a finite integer $m(a)\ge3$
such that its local section does not lift to $N/N^{(m(a))}$.
In particular its global section does not lift at every finite
derived length.  For each $h\ge1$ there is a uniform $m_h$ for
all $a\notin\ell^h\mathbb Z_\ell^{d_n}$.
\end{enumerate}
The assertion is existence of a finite obstruction, not a
calculation of the least $m(a)$ or of its obstruction cocycle.
It does not assert that one finite derived quotient has only
cuspidal sections, nor that the displayed family exhausts the
metabelian section set.

\paragraph{1. The arithmetic source is a quotient of the local group.}

Put $K=\mathbb Q(\zeta_\ell)$ and let $\mathfrak l$ be its
unique prime over $\ell$.  Write
\[
 P_{\rm ar}=G_{K,S_\ell}^{(\ell)},\qquad
 \Gamma_\ell=P_{\rm ar}\rtimes C_{\ell-1}
\]
for the arithmetic quotient used in D.20.  Let
$D\subset P_{\rm ar}$ be the image of $G_{K_{\mathfrak l}}$.
It is a closed subgroup.  We claim $D=P_{\rm ar}$.

Every continuous homomorphism $P_{\rm ar}\to\mathbb F_\ell$
which vanishes on $D$ defines a cyclic degree-$\ell$ extension
of $K$ unramified outside $\mathfrak l$ and split at
$\mathfrak l$.  Thus it is unramified at every finite prime.
The field $K$ is totally imaginary and has class number one,
as proved by the explicit Minkowski bound in D.20.  Global
class field theory implies that the character is zero.
Consequently the image of $D$ in the Frattini quotient
$P_{\rm ar}/\Phi(P_{\rm ar})$ is the whole quotient: otherwise
a nonzero linear functional annihilating that image would
give such a character.  A closed subgroup of a finitely
generated pro-$\ell$ group that surjects on the Frattini quotient
is the whole group.  For completeness, a proper closed subgroup
has proper image in some finite $\ell$-group quotient and lies
in the preimage of a maximal subgroup of that quotient; that
preimage is an index-$\ell$ subgroup.  This contradicts the
Frattini surjectivity.

Since $K_{\mathfrak l}/\mathbb Q_\ell$ has degree $\ell-1$,
the image of $G_{\mathbb Q_\ell}$ in $\Gamma_\ell$ also
surjects onto the constant quotient $C_{\ell-1}$.  It contains
the preceding image $P_{\rm ar}$ and hence equals all of
$\Gamma_\ell$.  Thus
\[
                 G_{\mathbb Q_\ell}\twoheadrightarrow\Gamma_\ell.
                                                        \tag{F.1}
\]
No assertion that the global absolute Galois group itself is
free is involved.

\paragraph{2. Every parameter remains nonzero locally.}

Recall $P=N^{(\ell)}$, $V=P^{\rm ab}$,
$B=\bigwedge^2 V$, and the $\Gamma_\ell$-equivariant embedding
$j:P'/P''\hookrightarrow N'/N''$ from D.20.  The parameter
module was chosen as a free submodule
\[
 W\simeq\mathbb Z_\ell^{d_n}
          \subset H^1(\Gamma_\ell,B).
\]
The constructed cocycle $s_a s_O^{-1}$ projects to the class
$a\in W$.  Inflation along the surjection (F.1) is injective
on continuous first cohomology.  This can be checked without
an inverse-limit assertion: if the inflated cocycle is a
coboundary with a vector $b\in B$, evaluating on a lift of
each element of $\Gamma_\ell$ gives exactly the same
coboundary formula on $\Gamma_\ell$.  Therefore every nonzero
parameter, and every nonzero difference of parameters, stays
nonzero in $H^1(G_{\mathbb Q_\ell},B)$.

To control geometric conjugators, set
$A=N^{\rm ab}$ and $k=\mathbb Q_\ell$.  We have
\[
 H^0(G_k,A)=0,\qquad H^0(G_k,A/I_O)=0.                 \tag{F.2}
\]
Here $I_O=\widehat{\mathbb Z}(1)$ is the primitive boundary
line at the rational cusp $O$.  Indeed, for each coefficient
prime $q$, compactified curve homology is $T_qJ$.
The group $J(k)$ has a torsion-free open subgroup of finite
index, obtained from the formal group sufficiently near zero;
hence its torsion subgroup is finite.  It follows that
$H^0(G_k,T_qJ)=T_q(J(k)_{\rm tors})=0$.  The boundary part
is a finite permutation module tensored with $\mathbb Z_q(1)$,
modulo its diagonal.  After a finite extension splitting the
cusps, both it and its quotient by the primitive rational line
$I_O$ are sums of $\mathbb Z_q(1)$.  These have no invariants:
a finite local field has only finitely many roots of unity,
so its $q$-adic cyclotomic image is infinite.  The homology
exact sequences give (F.2), prime by prime and hence for the
full Tate module.

If two local $s_a$ are conjugate, the conjugator's image in
$A$ is invariant because both abelian cocycles are zero.
By (F.2) the conjugator lies in $N'/N''$.  Projecting it first
to $P'/P''$ and then to $B$ would make the difference class a
coboundary.  The preceding inflation injectivity forces the
two parameters to be equal.  Thus geometric conjugation does
not destroy local injectivity.

\paragraph{3. The only geometric local class with zero abelian cocycle.}

The full local Kummer map for an abelian variety is injective:
$J(k)$ is a profinite group, so
$J(k)\to\varprojlim_m J(k)/mJ(k)$ is injective, and the Kummer
exact sequences identify this limit with its image in
$H^1(k,TJ)$.  If a local point or cusp section of $U_n$ has
zero abelian cocycle relative to the unit tangent at $O$, its
compactified point $x\in X_n(k)$ satisfies
$[x-O]=0$ in $J(k)$.  The Abel--Jacobi map of the smooth proper
curve $X_n$, of genus at least one, is injective on points:
a principal divisor $x-O$ for $x\ne O$ would give a degree-one
map to $\mathbb P^1$, contrary to positive genus.  Hence
$x=O$.  Since $O\notin U_n$, the section must lie in the cusp
packet at $O$.

The cohomology exact sequence of
$0\to I_O\to A\to A/I_O\to0$ and (F.2) show that
\[
                H^1(k,I_O)\longrightarrow H^1(k,A)
\]
is injective.  Therefore the only member of that packet with
zero abelian class is the unit tangential class $s_O$.
The metabelian centralizer/normalizer computation of D.20
shows that this is still the exact packet inside the full
metabelian quotient.  By Paragraph 2, $s_a$ is not that class
when $a\ne0$.  This proves the local nongeometric assertion.

D.54 now rules out a local lift of any $s_a$, $a\ne0$, to the
plain geometrically prosolvable group.  A global lift would
restrict to such a local lift, so it is ruled out as well.
In particular the plain five-power conclusion for this
specific family does not require a finite-Shafarevich theorem
for the Fermat quintic Jacobian.

\paragraph{4. Why the finite-derived conclusion is legitimate locally.}

We first prove the compactness fact that is unavailable over
$G_{\mathbb Q}$ with unrestricted ramification.  Let $k$ be a
finite extension of $\mathbb Q_p$.  For every finite solvable
$G_k$-group $Q$, the set $Z^1_{\rm cts}(G_k,Q)$ is finite.
For finite abelian $Q$, finiteness of $H^1(k,Q)$ is the local
finite-module theorem in Milne, \emph{Arithmetic Duality
Theorems}, Chapter I, Corollary 2.3 (and the finiteness proof
of Theorem 2.1), printed pages 27--28; each coboundary orbit has at
most $|Q|$ elements, so $Z^1$ is finite as well.  For general
finite solvable $Q$, use its last nonzero characteristic
derived subgroup $C$.  The target $H^1(k,Q/C)$ is finite by
induction.  The fibre over a represented class is an orbit
set of $H^1(k,C_\xi)$ for the corresponding twisted finite
abelian module.  The latter is finite by the same local
abelian theorem.  Thus $H^1(k,Q)$, and hence $Z^1(k,Q)$, is
finite.  This is finiteness of actual continuous cocycles,
not only a finiteness assertion about discrete abstract maps.

Let $M$ be any finitely generated prosolvable geometric group
with continuous $G_k$ action.  Its characteristic open
subgroups are cofinal.  Therefore
\[
 Z^1_{\rm cts}(G_k,M)
       =\varprojlim_V Z^1_{\rm cts}(G_k,M/V)             \tag{F.3}
\]
is compact Hausdorff: every finite coordinate is finite, and
a compatible family defines a continuous cocycle because
each finite coordinate is continuous.  Surjectivity of the
transition maps is not being assumed.  Quotient by geometric
conjugation is compact Hausdorff as well, since the action of
the compact group $M$ has closed orbits.

Put $M_m=N/N^{(m)}$, using the convention $M_2=N/N''$.
Suppose that a fixed actual metabelian local cocycle $s_a$
lifts up to geometric conjugation to $M_m$ for every $m$.
Conjugating a lift by an element lifting the required
metabelian conjugator makes its projection equal to the
\emph{actual} chosen cocycle.  Thus its exact lift fibre
$\mathcal L_m$ is a nonempty closed subset of the compact
space in (F.3).  The maps $M_{m+1}\to M_m$ give continuous
maps $\mathcal L_{m+1}\to\mathcal L_m$.  The inverse limit
of these nonempty compact spaces is nonempty: any finite
set of compatibility conditions is satisfied by choosing
an element at one sufficiently late stage and projecting
back, and compactness gives the full intersection.  This
uses no surjectivity of the fibre transition maps.

A compatible family defines an actual continuous cocycle
in
\[
        \varprojlim_m N/N^{(m)}=N^{\rm sol}
\]
and hence a full local prosolvable lift of $s_a$.  Paragraph 3
excludes this for $a\ne0$.  Therefore some finite stage
$m(a)\ge3$ has empty local lift fibre.  This is a genuine
finite-derived obstruction, not only a failure of an
arbitrarily chosen incompatible branch.

\paragraph{5. Uniformity away from the zero parameter.}

Let $L_m\subset\mathbb Z_\ell^{d_n}$ consist of parameters
whose local class lifts to $M_m$.  The local class map is
continuous, by the jointly continuous construction in D.20.
The image of the compact section space for $M_m$ in the
Hausdorff metabelian section space is closed; thus $L_m$
is closed.  The sets decrease with $m$, and the previous
paragraphs show
\[
                           \bigcap_{m\ge2}L_m=\{0\}.
\]
For fixed $h\ge1$, the set
$\mathbb Z_\ell^{d_n}\setminus\ell^h\mathbb Z_\ell^{d_n}$
is compact.  If it met every $L_m$, nested compactness would
give a nonzero element in their intersection.  Hence some
$m_h$ satisfies
\[
                         L_{m_h}\subset\ell^h\mathbb Z_\ell^{d_n}.
                                                        \tag{F.4}
\]
The same bound applies to global liftable parameters.

\paragraph{Remaining exact blocker.}

The proof does not identify $m(a)$ or $m_h$, nor a finite
solvable quotient and a numerical nonzero obstruction class.
Those are still the explicit obligations in Task D.5.
It also does not decide whether arbitrarily small nonzero
parameters survive each fixed stage, or whether other
noncuspidal families survive there.  Consequently the hunch
that no solvable quotient of fixed length has only cuspidal
sections is not decided by (F.4).
